% ============================================================
%  Section 4.X — Visual Grounding Evaluation (Major Comment 3)
%  Paste into Results / Experiments section of revised manuscript
% ============================================================

\subsection{Visual Grounding Evaluation}
\label{sec:visual_grounding}

Qualitative heatmaps alone are insufficient to support claims about model
grounding~\cite{chefer2021transformer}.
We therefore report a quantitative grounding protocol alongside revised
attribution visualizations.
Throughout, we refer to this analysis as \emph{visual grounding} rather than
clinical explainability: saliency maps indicate where the model attends when
producing an answer, but they are not validated by radiologist review.

\paragraph{Attribution maps.}
Saliency is computed from an ensemble of FineGrained text$\rightarrow$image
cross-attention maps (final three layers) and fusion-block
text$\rightarrow$image cross-attention, rather than raw ViT Grad-CAM.
This targets the multimodal fusion stage where question semantics interact
with grounded image tokens.

\paragraph{SLAKE (mask-based Pointing Game).}
SLAKE provides an official lesion segmentation mask (\texttt{mask.png}) for
each of the 642 test images.
We evaluate on the \textbf{full test set} ($n{=}1061$ QA pairs).
A pointing prediction is a \emph{hit} if the top 5\% saliency mass overlaps
the lesion mask; we also report argmax pointing for comparison.
Top-5\% overlap is more stable than argmax for small lesions.

Table~\ref{tab:pointing_game} summarizes the results.
MED-VLAT achieves 86.24\% VQA accuracy on SLAKE.
Mask-based Pointing Game accuracy is 43.92\% (argmax) and
\textbf{70.97\%} (top-5\% saliency overlap) on the same full test set.
On correctly answered questions only, top-5\% pointing remains 70.07\%,
indicating that grounding performance is not driven solely by incorrect
predictions.
Figure~\ref{fig:slake_grounding} shows five representative localization cases
on SLAKE with official \texttt{detection.json} bounding boxes (red) and
argmax pointing verification (column~e); the caption reports full-test
Pointing Game accuracy.

\paragraph{PathVQA and VQA-RAD (attention-focus proxy).}
PathVQA and VQA-RAD test annotations do not include segmentation masks or
bounding boxes, so mask-based Pointing Game is \textbf{not applicable} on either
dataset.
We therefore report \emph{mean attention focus}---the peak-to-mean saliency
ratio on the same combined attribution maps (higher values indicate more
localized, less diffuse saliency).
On PathVQA ($n{=}6719$), MED-VLAT achieves 65.66\% VQA accuracy (41.63\%
open, 89.65\% closed); mean attention focus is 2.21 on the full test set and
\textbf{1.99} on correctly answered questions, indicating tighter saliency when
predictions are right.
On VQA-RAD ($n{=}451$), the locked MED-VLAT dual-route system achieves
\textbf{81.82\%} test accuracy (369/451): FGGF for open and multi-choice
closed questions, dedicated closed yes/no specialist for binary questions.
VQA-RAD provides no region annotations, so Pointing Game is N/A; we report
mean attention focus \textbf{3.39} on localizing-style questions ($n{=}220$;
3.32 on the full test set) using FineGrained+fusion maps from the FGGF Q2A
checkpoint (see Table~\ref{tab:pointing_game}).
Without region labels, these scores serve as a complementary localization proxy
rather than verified grounding accuracy; SLAKE Pointing Game remains our
primary quantitative grounding metric.

\paragraph{Limitations.}
Transformer attribution methods remain an active research area and may
highlight high-frequency structures or background context.
A radiologist user study was not feasible within the revision timeline and is
identified as future work.

% --- Figure 6 (locked qualitative + Pointing Game layout) ---
\input{output/figures/slake_figure6_locked.tex}
